%% =====================================================================
%%  IEEE Signal Processing Letters (SPL) manuscript
%%  Spectral detection of backdoor triggers in graph neural networks
%%
%%  Class : IEEEtran.cls v1.8b (official IEEE class, supplied by the TeX distribution)
%%  Build : tectonic -X compile main.tex --outdir build
%%          or: pdflatex main && bibtex main && pdflatex main && pdflatex main
%%
%%  ---------------------------------------------------------------------
%%  Historical writing guide PAPER_MEMORY.md is not present in this checkout.
%%
%%  SOURCE MAP (which document owns which section):
%%    progress_report.tex -> background, motivation, attack taxonomy and
%%                           terminology, related work, citation pool
%%    final_report.tex    -> threat model, spectral representation,
%%                           clean-reference construction, detection
%%                           methodology, experimental setup, results
%%  Rewrite and compress; do not lift narrative from either report.
%%  Verify every reused citation actually supports the sentence it carries.
%%
%%  RESULTS PROVENANCE (every number in Section V):
%%    results/analysis/paper_graph_20260930
%%      statistical detector = pointwise_gaussian/holm (uncorrected rates
%%      in the text: pointwise_gaussian/none); MLP = supervised arm
%%      all__<attack>; GCN rows; mean over five folds.
%%      CA = clean_acc_before (poisoned model, clean test graphs);
%%      statistical-detector AUC = roc_auc of the holm row, whose score
%%      is the smallest Holm-adjusted p-value.
%%    TUDataset raw files (datasets/<name>/raw): fraction of graphs with
%%      n>100 is AIDS 0, MUTAG 0, NCI1 0.002, PROTEINS 0.060; Table III
%%      'Avg. deg.' is the mean over graphs of 2|E|/n.
%%    results/analysis/paper_graph_20260930_logreg
%%      logistic-regression ROC-AUC, same arm and rows.
%%  The pre-revision text is kept in main_before_revision.tex.
%%
%%  THE ARGUMENT (one chain, not a chronology):
%%    trigger insertion -> change in the normalized-Laplacian descriptor
%%    -> does the change exceed clean variation (statistical detector)
%%    -> does the whole descriptor separate individual graphs (MLP)
%%    -> does rejecting flagged graphs suppress the backdoor -> limits
%%
%%  Intermediate detectors from the reports are development history and
%%  supporting evidence, not competing proposed methods. The MLP-based
%%  spectral detector is the proposed method.
%%  ---------------------------------------------------------------------
%%
%%  SPL length rules: 4 pages of technical content; a 5th page is allowed
%%  for References only. Author biographies/photographs are not permitted.
%%  Do NOT change margins, font sizes, column widths or line spacing.
%% =====================================================================
\RequirePackage[T1]{fontenc}
\documentclass[journal]{IEEEtran}

% ---------------------------------------------------------------------
% Packages -- keep this list minimal; add only what the paper really uses
% ---------------------------------------------------------------------
\usepackage{cite}                       % IEEE-style [1]-[3] range compression
\usepackage{amsmath,amssymb,amsfonts}   % equations and math symbols
\usepackage{graphicx}                   % figures
\usepackage{booktabs}                   % \toprule / \midrule / \bottomrule
\usepackage{tabularx}                   % Y columns in the configuration table
\usepackage{url}
\usepackage{xcolor}
\usepackage[hidelinks]{hyperref}        % keep last

\graphicspath{{figures/}}

% ---------------------------------------------------------------------
% Table helpers
% ---------------------------------------------------------------------
\newcolumntype{Y}{>{\raggedright\arraybackslash}X}
\newcommand{\tabstyle}{%
  \footnotesize
  \setlength{\tabcolsep}{3pt}%
  \renewcommand{\arraystretch}{1.15}%
}

% ---------------------------------------------------------------------
% Notation macros -- one symbol per quantity
% ---------------------------------------------------------------------
\newcommand{\bA}{\mathbf{A}}            % adjacency matrix
\newcommand{\bD}{\mathbf{D}}            % degree matrix
\newcommand{\bX}{\mathbf{X}}            % node feature matrix
\newcommand{\bL}{\mathbf{L}}            % normalized Laplacian
\newcommand{\bs}{\mathbf{s}}            % spectral descriptor
\newcommand{\calV}{\mathcal{V}}         % node set
\newcommand{\calE}{\mathcal{E}}         % edge set
\newcommand{\calD}{\mathcal{D}}         % dataset
\newcommand{\calL}{\mathcal{L}}         % training loss
\newcommand{\fth}{f_{\theta}}           % GNN model with parameters theta

\newtheorem{definition}{Definition}

\begin{document}

% Note on wording: "attack-agnostic", used in the progress report, is
% deliberately avoided -- the learned detector is trained and tested on the
% same attack, and cross-attack generalization is not addressed here.
\title{On the Imperceptibility of Graph Backdoor Trojans\\ in Spectral Domain}

% TODO: authors, IEEE membership grades, affiliations, funding note.
% Check the SPL submission policy for the current review cycle before
% filling this in (anonymized vs. non-anonymized submission).
\author{TODO: Author Names%
\thanks{TODO: submission date, funding and sponsor acknowledgment.}%
\thanks{TODO: author affiliations, addresses and e-mail addresses.}}

% Running head; short author form on the right, e.g. "Author \MakeLowercase{\textit{et al.}}: Short Title"
\markboth{IEEE Signal Processing Letters}%
{TODO: Author \MakeLowercase{\textit{et al.}}: TODO Short Title}

\maketitle

\begin{abstract}

\end{abstract}

\begin{IEEEkeywords}
Backdoor attacks, graph neural networks, graph signal processing, normalized
Laplacian, spectral detection.
\end{IEEEkeywords}

\IEEEpeerreviewmaketitle

% =====================================================================
\section{Introduction}
\IEEEPARstart{G}{raph} neural networks~\cite{gcn,gin} are widely used to classify
graph-structured data, but they can be compromised by backdoor attacks
introduced through poisoned training data~\cite{sba2021,gta2021}. In graph
classification, an attacker embeds a trigger in selected training graphs and
relabels them with an attacker-chosen target class. The trained model then
behaves normally on clean graphs but assigns triggered graphs to the target
class. Consequently, clean-test accuracy alone cannot establish the absence of
backdoor behavior.

Many attacks try to keep triggers imperceptible: the attacker limits how much a
trigger may alter a graph, so that poisoned graphs resemble clean ones
(Section~\ref{sec:policy}), and some attacks, such as
DPSBA~\cite{dpsba2025}, constrain the structural and semantic deviations of
poisoned graphs explicitly. This letter asks whether structural triggers remain
imperceptible in the spectral domain, that is, whether the normalized-Laplacian
spectrum of a triggered graph can be told apart from those of clean graphs;
detectability is thus the converse of imperceptibility. Inserting a trigger
rewires edges, which can move the eigenvalues of the normalized graph
Laplacian~\cite{chung1997spectral}. The question is whether these changes are
consistent enough to distinguish individual triggered graphs from clean ones.

The evaluation proceeds in three steps. First, a statistical test compares the
spectrum of each graph with the variation among clean graphs, to determine
whether trigger-induced changes exceed this natural variation. Second, a
multilayer perceptron (MLP) trained on labelled clean and triggered spectra
tests whether the spectrum as a whole supports individual-graph detection.
Third, graphs flagged by either detector are rejected before they reach the
target model, and the remaining attack success is measured. The MLP threshold
is set on held-out clean graphs so that the two detectors operate at
approximately the same cost on clean graphs.

The contributions are: (i) an assessment, through a clean-reference test, of
whether the normalized-Laplacian changes induced by structural backdoor
triggers exceed the natural spectral variation among clean graphs; (ii) a
spectrum-only MLP detector for individual graphs that requires no access to
the target model at inference; and (iii) an evaluation of how rejecting
detected graphs affects attack success while retaining clean inputs.

% =====================================================================
\section{Related Work}
Structural graph backdoors differ in how triggers are constructed and
inserted. The subgraph-based backdoor attack (SBA) uses a reusable
random-graph template~\cite{sba2021}, while Motif-Backdoor employs motif-based
templates with model-guided placement~\cite{motif}. The placement
strategy may depend on the host graph even when the trigger structure
itself remains fixed. In contrast, the graph Trojan attack (GTA) constructs
graph-specific triggers using a learned generator~\cite{gta2021}, while TRAP uses
gradient-guided edge perturbations~\cite{10.1145/3545948.3545976}.

Backdoor attacks also differ in the graph components they perturb.
DPSBA constrains structural and semantic deviations through
distribution-aware clean-label poisoning~\cite{dpsba2025}. Other
attacks operate primarily on node features: SPEAR preserves graph
topology~\cite{spear2025}, while spectral-domain attacks embed triggers
in the spectral coefficients of node-feature
signals~\cite{spectralba2024}. Because such feature-only attacks leave
the adjacency matrix unchanged, they also leave the normalized
Laplacian spectrum unchanged. They therefore fall outside the scope of
the structural detector considered here.

Existing graph backdoor defenses differ substantially in their access
requirements. GSP purifies training graphs using spectral low-pass
filtering~\cite{gsp2025}, while RIGBD~\cite{rigbd2025} and
DShield~\cite{dshield2025} inspect potentially poisoned training data
and intervene during training. DMGNN relies on model-dependent
counterfactual explanations for trigger pruning~\cite{dmgnn2024},
whereas GraphProt avoids access to model internals but requires
repeated classifier queries over sampled subgraphs~\cite{graphprot2025}.
The SSP working paper combines feature-covariance eigenanalysis with
neighborhood feature deviations~\cite{ssp2026}, rather than the adjacency
Laplacian eigenvalues studied here. A2GBD uses active node selection and
reinforcement learning for attack-agnostic defense~\cite{a2gbd2026}.
The proposed detector operates directly on the normalized
Laplacian spectrum of an incoming graph. Once trained on labelled clean
and triggered spectra, the MLP requires neither the target's training
set nor access to its parameters, gradients, or internal
representations, and it does not require querying the target for
detection. The detector therefore acts as a pre-inference filter rather
than modifying the input graph or retraining the target.

% =====================================================================
\section{Preliminaries and Problem Formulation}

\subsection{Graph Classification}
A graph is $G=(\calV,\calE,\bX)$ with $n=|\calV|$ nodes, adjacency matrix
$\bA\in\{0,1\}^{n\times n}$, degree matrix $\bD$ and node features
$\bX$. A graph classifier $\fth$ with parameters $\theta$ maps $G$ to a vector
of class probabilities.
\subsection{Holm Procedure}
Multiplicity control throughout this letter uses the Holm step-down
procedure~\cite{holm1979}; each family of hypotheses consists of the
coordinatewise tests of a single graph (Section~\ref{sec:stat}).
\begin{definition}[Holm Step-Down Procedure]
Consider $m$ null hypotheses $H_1,\dots,H_m$ with corresponding valid
$p$-values $p_1,\dots,p_m$. Let
\[
  p_{(1)}\leq p_{(2)}\leq\dots\leq p_{(m)}
\]
denote the ordered $p$-values, and let $H_{(i)}$ be the hypothesis associated
with $p_{(i)}$. For a prescribed family-wise significance level
$\alpha\in(0,1)$, define
\begin{equation}
  \ell^{*}=\max\Bigl(\{0\}\cup\bigl\{\ell\in\{1,\dots,m\}:\,
  p_{(i)}\leq\tfrac{\alpha}{m-i+1}\ \,\forall i\leq \ell\bigr\}\Bigr).
  \label{eq:holm-step-down}
\end{equation}
The procedure rejects $H_{(1)},\dots,H_{(\ell^{*})}$. Thus, testing starts with
the smallest $p$-value and stops at the first failed comparison; that
hypothesis and all subsequent hypotheses are not rejected.

Equivalently, the Holm-adjusted $p$-value for $H_{(i)}$ is
\begin{equation}
  p^{\mathrm{Holm}}_{(i)}
  =\min\left\{
    1,\;
    \max_{1\leq t\leq i}
    \left[(m-t+1)\,p_{(t)}\right]
  \right\}.
  \label{eq:holm-adjusted}
\end{equation}
A hypothesis is rejected when its adjusted $p$-value does not exceed $\alpha$.
For valid input $p$-values, this procedure controls the probability of
rejecting at least one true null hypothesis at level $\alpha$, without
requiring independence among the tests.
\end{definition}
\subsection{Backdoor Trojan}
\subsubsection{Basics of a Backdoor Trojan }
A backdoor Trojan (backdoor, for short) is a hidden behavior implanted in a
model through poisoned training data and activated by a trigger, a pattern
the attacker inserts into an input graph.

\begin{definition}[Backdoor Attack]
Let $\mathcal{X}$ be the space of attributed graphs, let $\mathcal{C}$ be a
finite label set with $C=|\mathcal{C}|$, and let $\mathcal{P}$ be a
distribution over $\mathcal{X} \times \mathcal{C}$. Let
$f_\theta : \mathcal{X} \to \Delta^{C-1}$ be a model trained on a sample from
$\mathcal{P}$, and write
$h_\theta(G)=\arg\max_{c\in\mathcal{C}} f_{\theta,c}(G)$ for its predicted
label. Let $y_t \in \mathcal{C}$ be a fixed target label and let
$T : \mathcal{X} \times \mathcal{C} \to \mathcal{X}$ be a trigger policy, so
that $T(G,y_t)$ is the triggered version of $G$. For tolerances
$\delta_a,\delta_c \in [0,1)$, a backdoor attack produces a poisoned model
$f_{\theta^*}$ such that
\begin{align*}
\Pr_{(G,y)\sim\mathcal{P}}\bigl[\,h_{\theta^*}\bigl(T(G,y_t)\bigr)=y_t
  \;\big|\; y\neq y_t \,\bigr] &\ \ge\ 1-\delta_a,\\
\Pr_{(G,y)\sim\mathcal{P}}\bigl[\,h_{\theta^*}(G)=h_{\theta}(G)\,\bigr]
  &\ \ge\ 1-\delta_c .
\end{align*}
The first condition is the backdoor behaviour on triggered inputs and the
second is preservation of clean behaviour; $\delta_a=\delta_c=0$ recovers the
idealized deterministic case.
\end{definition}

\begin{definition}[Clean Accuracy]
Let $\mathcal{D}_{\text{test}}$ be a held-out test set, disjoint from the
training set $\mathcal{D}$, and let $\Pr_{(G,y)\sim\mathcal{D}_{\text{test}}}$
denote the empirical probability over it. The accuracy of $f_{\theta^*}$ on
clean, untriggered inputs is
\[
\mathrm{CA} = \Pr_{(G,y) \sim \mathcal{D}_{\text{test}}}\!\left[h_{\theta^*}(G) = y\right].
\]
\end{definition}

\begin{definition}[Attack Success Rate]
The rate at which triggered inputs are mapped to the target label $y_t$:
\[
\mathrm{ASR} = \Pr_{(G,y) \sim \mathcal{D}_{\text{test}}}\!\left[h_{\theta^*}(T(G,y_t)) = y_t \;\middle|\; y\neq y_t\right].
\]
\end{definition}

\subsubsection{Backdoor Insertion Framework}\label{sec:policy}

The common backdoor attack framework consists of three stages~\cite{sba2021,gta2021}: (1)
trigger-insertion policy design, where the attacker constructs or learns $T$;
(2) data poisoning, where $T$ is applied to a subset of the training data to
build a poisoned dataset $\widetilde{\mathcal{D}}$; and (3) model poisoning,
where a (possibly third-party) model is trained on $\widetilde{\mathcal{D}}$,
implanting the backdoor.

Let $f_\theta : \mathcal{X} \to \Delta^{C-1}$ be the target model and
$\mathcal{D} = \{(G_i, y_i)\}_{i=1}^{N}$, drawn i.i.d.\ from $\mathcal{P}$,
the clean training set.

\paragraph{Trigger-Insertion Policy Design}
The attacker defines a policy $T : \mathcal{X} \times \mathcal{C} \to
\mathcal{X}$. $T$ may be non-parametric (e.g., solved per
instance) or restricted to a parametric family
$\{T_\phi\}_{\phi\in\mathcal{W}}$, i.e.\ the collection of policies obtained as the
generator weights $\phi$ range over the generator weight space $\mathcal{W}$, learned once
and reused at inference. Given a surrogate $f_{\hat\theta}$,
\begin{align}
T^* = \arg\max_{T} \; & \frac{1}{N}\sum_{i=1}^{N} \mathbf{1}\bigl[h_{\hat\theta}(T(G_i,y_t))=y_t\bigr] \notag \\
\text{s.t.} \; & d\big(T(G_i,y_t), G_i\big) \le \epsilon, \quad i=1,\dots,N,
\end{align}
where $d$ is a distortion measure and $\epsilon$ bounds perceptibility.

For $T = T_\phi$, i.e. $T$ is realized by a neural network with weights $\phi
\in \mathcal{W}$, the search over policies reduces to optimizing $\phi$:
\begin{align}
\phi^* = \arg\max_{\phi \in \mathcal{W}} \; & \frac{1}{N}\sum_{i=1}^{N} \mathbf{1}\bigl[h_{\hat\theta}(T_\phi(G_i,y_t))=y_t\bigr] \notag \\
& - \; \frac{\gamma}{N}\sum_{i=1}^{N} d\bigl(T_\phi(G_i,y_t), G_i\bigr),
\label{eq:phi-star}
\end{align}
with the averaged distortion as a soft penalty replacing the hard constraint.
The indicator above is not differentiable in $\phi$, so it is relaxed to the
cross-entropy surrogate $\mathcal{L}_{\mathrm{CE}}(\cdot,y_t)$, which is small
exactly when $f_{\hat\theta}$ concentrates its mass on $y_t$. Since
$f_{\hat\theta}$ and $T_\phi$ are differentiable, $\phi$ is then trained by
gradient descent on this surrogate:
\begin{align}
\mathcal{J}(\phi) &= \frac{1}{N}\sum_{i=1}^{N} \Big[ \mathcal{L}_{\mathrm{CE}}\big(f_{\hat\theta}(T_\phi(G_i,y_t)), y_t\big) \notag \\
&\qquad\qquad + \gamma \, d\big(T_\phi(G_i,y_t), G_i\big) \Big], \\
\phi &\leftarrow \phi - \eta \nabla_\phi \mathcal{J}(\phi).
\end{align}
This amortizes trigger construction: $\phi^*$ is learned once over
$\mathcal{D}$, and $T_{\phi^*}(G, y_t)$ is then generated in one forward pass
for any new $(G, y_t)$, rather than solving a fresh optimization per pair.

\paragraph{Data Poisoning}
For poisoning rate $\rho \in (0,1)$, $\mathcal{D}_p \subset \mathcal{D}$,
$|\mathcal{D}_p| = \lfloor\rho N\rfloor$,
\begin{align}
\mathcal{D}_p' &= \big\{ (T^*(G, y_t), y_t) \mid (G,y) \in \mathcal{D}_p \big\}, \\
\widetilde{\mathcal{D}} &= (\mathcal{D} \setminus \mathcal{D}_p) \cup \mathcal{D}_p'.
\end{align}
The attacker aims to preserve clean accuracy while increasing attack success;
this tradeoff is evaluated empirically.

\paragraph{Model Poisoning}
The target model's owner trains on $\widetilde{\mathcal{D}}$ as released,
without attacker control:
\begin{equation}
\theta^* = \arg\min_{\theta} \; \frac{1}{|\widetilde{\mathcal{D}}|} \sum_{(G,y) \in \widetilde{\mathcal{D}}} \mathcal{L}\big(f_\theta(G), y\big),
\end{equation}
where $\mathcal{L}$ is the training loss.

\subsection{Trigger Insertion in Graphs}\label{sec:insertion}
A trigger can either inject new nodes or modify existing host nodes. Let
$\calV_g$ be the placed trigger nodes, $\calE_g$ the inserted edges (including any
links to the host), and $\calE_{\mathrm{cut}}\subseteq\calE$ the removed edges.
Writing $\mathbf{x}^{g}_v$ for the assigned feature row at a trigger node,
the general composition is
\begin{align}
\tilde{\calV} &= \calV\cup\calV_g, \label{eq:subgraph-node-composition}\\
\tilde{\calE} &= (\calE\setminus\calE_{\mathrm{cut}})\cup\calE_g,
\label{eq:subgraph-edge-composition}\\
\tilde{\mathbf{x}}_v &=
\begin{cases}
\mathbf{x}^{g}_v, &v\in\calV_g,\\
\mathbf{x}_v, &v\in\calV\setminus\calV_g.
\end{cases}\label{eq:subgraph-feature-composition}
\end{align}
Consequently,
\begin{equation}
|\tilde{\calV}|=|\calV|+|\calV_g\setminus\calV|.
\label{eq:subgraph-node-count}
\end{equation}
Node-injection attacks use $\calV_g\cap\calV=\varnothing$, as in
UGBA~\cite{ugba2022} and DPGBA~\cite{dpgba2023}. In contrast, all three
graph-classification attacks evaluated here select existing nodes,
$\calV_g\subseteq\calV$, and preserve the vertex count. Their nominal trigger
size $k=|\calV_g|$ is given in Table~\ref{tab:setup}.

\subsubsection{Static Triggers}
\label{sec:static}
A static trigger uses a fixed pattern $\Gamma=(\calV_\Gamma,\calE_\Gamma)$ and a
placement map $\iota_\pi:\calV_\Gamma\to\calV$ selected by a policy $\pi$. For
the SBA and Motif implementations evaluated here,
$\calV_g=\iota_\pi(\calV_\Gamma)$ and
\begin{equation}
T(G,y_t)=\bigl(\calV,\,(\calE\setminus\calE[\calV_g])\cup\iota_\pi(\calE_\Gamma),\,\bX\bigr),
\label{eq:static-trigger}
\end{equation}
where $\calE[\calV_g]$ contains the edges with both endpoints in $\calV_g$, and
$\iota_\pi(\calE_\Gamma)$ maps pattern edges to those hosts. Edges within the
selected region are replaced; external connections and features are preserved.
In the notation of~\eqref{eq:subgraph-node-composition}--\eqref{eq:subgraph-feature-composition},
this sets $\calE_{\mathrm{cut}}=\calE[\calV_g]$, $\calE_g=\iota_\pi(\calE_\Gamma)$,
and $\mathbf{x}^{g}_v=\mathbf{x}_v$.

SBA uses an Erd\H{o}s--R\'enyi pattern~\cite{sba2021}; our implementation
places it on uniformly sampled host nodes. Motif-Backdoor uses motif
structure with model-guided placement~\cite{motif}. Here the motif is M44, a
four-node triangle with a tail, and high-degree candidates are ranked by the
change in surrogate logits when their incident edges are removed. The pattern is
fixed, but its spectral effect can vary with host topology and placement.

\subsubsection{Dynamic Triggers}
\label{sec:dynamic}
A dynamic policy $T_\phi$ generates perturbations conditioned on the host.
The evaluated GTA implementation selects a region $\calV_g\subseteq\calV$ and
uses topology and feature generators~\cite{gta2021}:
\begin{equation}
\begin{aligned}
\tilde{\bA}&=\max\bigl(\bA,\,\bA_\phi(G)\bigr),\\
\tilde{\bX}&=\bX+\bX_\phi(G).
\end{aligned}
\label{eq:dynamic-trigger}
\end{equation}
The generated adjacency matrix $\bA_\phi(G)$ is binarized and masked to pairs
inside $\calV_g$, and the maximum is taken entrywise; the feature perturbation
$\bX_\phi(G)$ is masked to the rows of $\calV_g$. Thus GTA adds edges within
the selected region and changes features, while keeping $\tilde{\calV}=\calV$.
In the notation of~\eqref{eq:subgraph-node-composition}--\eqref{eq:subgraph-feature-composition},
$\calE_{\mathrm{cut}}=\varnothing$, $\calE_g$ holds the generated edges, and
$\mathbf{x}^{g}_v$ is the perturbed feature row of a node of $\calV_g$. The
region selection itself is not learned. The general optimization framework in
Section~\ref{sec:policy} is a surrogate objective, not a guarantee that a
hard distortion bound is satisfied. Related imperceptibility objectives
include distribution-aware poisoning in DPSBA~\cite{dpsba2025} and
feature-only perturbations in SPEAR~\cite{spear2025}; these are distinct
attack mechanisms.

\subsection{Graphs in Spectral Domain}\label{sec:spectral}

\begin{definition}[Normalized Laplacian Spectrum]
For a graph $G=(\calV,\calE,\bX)$ with adjacency matrix $\bA$ and degree matrix
$\bD$, the normalized Laplacian is
\[
\bL = \mathbf{I} - \bD^{-1/2}\bA\bD^{-1/2}.
\]
For isolated nodes, the corresponding entry of $\bD^{-1/2}$ is
set to zero, so each contributes an eigenvalue of $1$. The eigenvalues of
$\bL$,
\[
\boldsymbol{\lambda}(G)=[\lambda_1,\dots,\lambda_n]^{\top},
\qquad
0\le\lambda_1\le\lambda_2\le\dots\le\lambda_n\le 2,
\]
are referred to as the spectrum of $G$. The spectrum is invariant to node
ordering and bounded independently of graph size~\cite{chung1997spectral},
which is what allows graphs to be compared without aligning their nodes.
\end{definition}

Following graph signal processing~\cite{shuman2013gsp}, the eigenvectors of
$\bL$ serve as the oscillation modes of the graph and the eigenvalues as their
frequencies. An eigenvector assigns a value to every node, and its eigenvalue
measures how strongly these values, scaled by the square root of the node
degree, differ across the edges of the graph~\cite{chung1997spectral}.

Low frequencies, the small eigenvalues, belong to eigenvectors that change
slowly across edges, so that connected nodes receive similar values. The
number of zero eigenvalues equals the number of connected components that
contain at least one edge, because each such component carries one
eigenvector that is perfectly smooth on it. For a connected graph, a small
nonzero $\lambda_2$ indicates a bottleneck: the graph can be split into two
sizeable parts by removing comparatively few edges~\cite{chung1997spectral}.
High frequencies, the large eigenvalues, belong to eigenvectors whose values
alternate between neighboring nodes. The largest possible value, $2$, occurs
exactly when a connected component with an edge is bipartite, that is, when
its nodes can be divided into two groups such that every edge joins the two
groups~\cite{chung1997spectral}.

The spectrum does not determine a graph uniquely, but it summarizes how the
graph is connected: how many separate pieces it has, whether it has a
bottleneck, and whether it is bipartite. A trigger rewires edges among a few
nodes (Section~\ref{sec:insertion}), which may change this connectivity and
with it some eigenvalues.

The spectrum of a graph on $n$ vertices is a vector of length $n$, so the
spectra of two graphs of different order live in spaces of different
dimension and cannot be compared coordinatewise. Two devices remove this
dependence on size: the eigenvalues are reindexed by their \emph{relative}
position in the spectrum, and the resulting profile is resampled on a grid
that is shared by all graphs.

\begin{definition}[Normalized Spectral Profile]
Let $G$ be a graph on $n\ge 2$ vertices with spectrum $\boldsymbol{\lambda}(G)$.
Assign to each eigenvalue the normalized position
\begin{equation}
r_i=\frac{i-1}{n-1},\qquad i=1,\dots,n,
\label{eq:rindex}
\end{equation}
so that $r_1=0$ and $r_n=1$ irrespective of $n$. The \emph{spectral
profile} of $G$ is the unique piecewise-linear function
$\Lambda_G\colon[0,1]\to[0,2]$ with breakpoints $r_1,\dots,r_n$ satisfying
$\Lambda_G(r_i)=\lambda_i$.
\end{definition}

Between two breakpoints, the profile interpolates linearly between
consecutive eigenvalues, so it is nondecreasing and takes values in $[0,2]$.
In effect, $\Lambda_G$ retains the \emph{shape} of the spectrum while
discarding its length: the horizontal axis measures how far one has advanced
through the spectrum rather than how many eigenvalues there are.

Graphs of different order are compared on a shared grid of $K$ equally spaced
positions $j/K$, $j=1,\dots,K$. The position $0$ is omitted, since
$\Lambda_G(0)=\lambda_1$ equals zero for every graph with an edge and
therefore carries no information. Reading the profile on this grid gives the
spectral descriptor
\begin{equation}
\bs(G)=\Bigl[\Lambda_G\bigl(1/K\bigr),\,
\Lambda_G\bigl(2/K\bigr),\,\dots,\,
\Lambda_G\bigl(1\bigr)\Bigr]^{\!\top}\!\in[0,2]^{K},
\label{eq:specvec}
\end{equation}
whose $j$th coordinate is denoted $s_j(G)$ and whose last entry is
$\Lambda_G(1)=\lambda_n$. The map $G\mapsto\bs(G)$ sends every graph on at
least two vertices into the same $K$-dimensional space, so graphs of different
order become directly comparable. It inherits permutation invariance from the
spectrum, its coordinates are nondecreasing by construction, and they are
confined to $[0,2]$, which bounds the input range. For $n-1>K$ the grid
subsamples a dense spectrum, while for $n-1<K$ it interpolates between
consecutive eigenvalues. In both cases, the first coordinates describe the
low-frequency, connectivity end of the spectrum, and the last coordinates the
high-frequency, bipartite-like end near $\lambda_n$.

% =====================================================================
\section{Fingerprint of Backdoor Trojan}
We call the change that a trigger makes to the descriptor~\eqref{eq:specvec}
its spectral fingerprint. Since the spectrum depends only on the edges, the
fingerprint reflects the rewired edges of Section~\ref{sec:insertion}; feature
perturbations, such as those of GTA, do not affect it. At inference time the
defender observes a single graph, not its clean version, so the fingerprint
cannot be measured as a difference. The question is instead whether the
descriptor of an observed graph lies outside the variation of clean
descriptors. Two detectors address this question. The statistical detector of
Section~\ref{sec:stat} needs only clean graphs; it tests whether
trigger-induced deviations exceed natural clean variation and serves as the
reference. The MLP detector of Section~\ref{sec:mlp} is the proposed method; it
learns a joint decision over all coordinates from labelled clean and triggered
graphs.

\subsection{Statistical Detection of Backdoor Trojan}\label{sec:stat}
The statistical detector compares the descriptor of an observed graph with a
reference built from clean graphs alone.

\begin{definition}[Clean-Reference Statistical Detector]
Let $\mathcal{R}$ be a reference set of clean graphs, and let $\mu_j$ and
$\sigma_j$ be the sample mean and standard deviation of the $j$th coordinate
$s_j$ of the descriptor over $\mathcal{R}$. For an observed graph $G$, the
standardized deviation at coordinate $j$ and its two-sided Gaussian $p$-value
are
\begin{equation}
z_j(G)=\frac{s_j(G)-\mu_j}{\sigma_j},
\qquad
p_j(G)=2\bigl[1-\Phi\bigl(|z_j(G)|\bigr)\bigr],
\label{eq:clean-reference-test}
\end{equation}
where $\Phi$ is the standard normal cumulative distribution function. The
detector applies the Holm procedure to the $p$-values of $G$ at level $\alpha$
and flags $G$ if at least one coordinate is rejected:
\begin{equation}
\hat{b}_{\mathrm{stat}}(G)=
\begin{cases}
1, & \text{Holm rejects at least one } p_j(G),\\
0, & \text{otherwise}.
\end{cases}
\label{eq:statistical-detector}
\end{equation}
\end{definition}

\noindent
Triggered graphs never enter $\mathcal{R}$. Since the decision asks only
whether any coordinate is rejected, it is equivalent to comparing the smallest
$p_j(G)$ with $\alpha/m$, where $m$ is the number of tested coordinates; the
later steps of the Holm procedure do not change it. Because the Gaussian model
is an approximation, $\alpha$ does not fix the rate at which clean graphs are
flagged; this rate is measured on clean test graphs.

\subsection{Model-Based Backdoor Detection}\label{sec:mlp}
The statistical detector tests each coordinate separately and ignores
dependencies between coordinates. The model-based detector instead learns
one decision over the whole descriptor.

\begin{definition}[MLP-Based Spectral Detector]
Let the training set consist of clean graphs and their triggered versions,
labelled $b=0$ and $b=1$, respectively, and let $\bar{\bs}(G)$ denote the
descriptor with each coordinate standardized by statistics of the training
graphs. An MLP $g_\psi$ with parameters $\psi$, trained with the binary
cross-entropy loss, maps it to the detection score
$q_\psi(G)=g_\psi\bigl(\bar{\bs}(G)\bigr)$. For a clean-rejection budget
$\beta$, a threshold $\tau$ is set on held-out clean graphs so that the
fraction of them flagged does not exceed $\beta$. The detector flags $G$ when
its score exceeds the threshold:
\begin{equation}
\hat{b}_{\mathrm{MLP}}(G)=
\begin{cases}
1, & q_{\psi}(G)>\tau,\\
0, & \text{otherwise}.
\end{cases}
\label{eq:mlp-detector}
\end{equation}
\end{definition}

\noindent
Logistic regression on the same standardized descriptor serves as a linear
reference.

Training requires triggered examples of the attack being detected. A separate
detector is trained for each dataset and attack, and how a detector behaves on
attacks it was not trained on is not addressed in this letter. Because a
defender does not normally hold such examples, the MLP results indicate what a
supervised detector can achieve on the descriptor; they do not describe a
defense against unseen attacks.

% =====================================================================
\section{Numerical Results}

\subsection{Simulation Setup}
\label{sec:setup}
Table~\ref{tab:setup} lists the configuration. The target model is a
two-layer graph convolutional network (GCN)~\cite{gcn}. It is trained on five stratified folds, each holding
out $20\%$ of the graphs for testing and $10\%$ for validation, and is
evaluated at its best-validation checkpoint (Appendix~\ref{app:training}).

\begin{table}[t]
\caption{Experimental configuration.}
\label{tab:setup}
\centering
\tabstyle
\begin{tabularx}{\columnwidth}{@{}lY@{}}
\toprule
Component & Setting \\
\midrule
Target model      & two-layer GCN, hidden dimension $64$, dropout $0.5$ \\
Target training   & Adam, lr $10^{-3}$, weight decay $5\times10^{-4}$, batch $32$, up to $50$ epochs, early stopping patience $30$ \\
Target split      & $70/10/20$ train/validation/test, five stratified folds \\
Attack budget     & target class $y_t=0$; poisoning rate $\rho=0.10$; trigger size $k=4$ \\
Descriptor        & normalized Laplacian; $K=100$ grid positions; positions with clean standard deviation below $10^{-6}$ are not used \\
Statistical detector & two-sided Gaussian $p$-values; Holm at $\alpha=0.05$; $10^{-6}$ added to $\sigma_j$ \\
Detector pool split & $70\%$ of source graphs for fitting, $30\%$ for threshold calibration \\
Supervised detectors & logistic regression; MLP with one hidden layer of $64$ ReLU units \\
MLP optimization  & Adam, lr $10^{-3}$, $200$ epochs, batch $256$ \\
Budget $\beta$    & clean rejection rate of the statistical detector on the pool, leave-one-out \\
\bottomrule
\end{tabularx}
\end{table}

The detectors of each fold are built from the graphs outside its test split,
called the pool, so no test graph influences a detector. The reference set
$\mathcal{R}$ of the statistical detector contains all clean pool graphs. For
the MLP and the logistic regression, the pool is split by source graph, so
that a clean graph and its triggered version always fall on the same side.
The fitting part provides the training pairs; the clean graphs of the
calibration part set the threshold $\tau$. The budget $\beta$ is the rate at
which the statistical detector rejects clean pool graphs, each scored against
a reference that excludes it. Both detectors are therefore compared at
approximately the same cost on clean graphs.

% graphs with n<=100 (TUDataset raw files): AIDS 100%, MUTAG 100%, NCI1 99.83%, PROTEINS 93.98%
With $K=100$, the grid interpolates between eigenvalues for at least $93\%$ of
the graphs of every dataset ($n\le K$), so neighboring coordinates of the
descriptor are strongly dependent. Holm's procedure does not require
independence, and the influence of $K$ is not studied in this letter.

All metrics are computed on the test split and averaged over the five folds.
The clean accuracy (CA) and the attack success rate (ASR) are those of the
poisoned target model. The false-positive rate (FPR) is the fraction of clean
test graphs flagged. The true-positive rate (TPR) is the fraction of triggered
test graphs flagged, taken over the eligible graphs: those whose true label
differs from $y_t$ and to which the trigger can be applied. The ASR is computed
over the same eligible graphs, which adds the second requirement to the
conditioning event in its definition. The area under the receiver operating
characteristic curve (ROC-AUC) summarizes how well the detector score separates
clean and triggered graphs over all thresholds; graphs are ranked by $q_\psi$
for the MLP, by the analogous output for logistic regression, and by the
smallest Holm-adjusted $p$-value for the statistical detector. For
the defense, flagged graphs are rejected before they reach the target model.
The post-detection ASR counts a triggered graph as a successful attack only if
it is accepted and classified as $y_t$; rejected graphs count as failed
attacks. It is measured by running the target model on the accepted graphs,
not derived from the TPR, because detection and attack success need not occur
on the same graphs.

\subsubsection{Datasets}
AIDS, MUTAG, and NCI1 contain small molecules and are called the molecular
datasets below. PROTEINS contains protein structures. All four come from the
TUDataset benchmark~\cite{tudataset}; Appendix~\ref{app:structure} summarizes
their structure (Table~\ref{tab:dataset_stats_graph}).

\subsubsection{Backdoor Trojan Attacks}
Each dataset is poisoned with SBA and Motif-Backdoor (static triggers,
Section~\ref{sec:static}) and with GTA (dynamic trigger,
Section~\ref{sec:dynamic}), at the poisoning rate and trigger size of
Table~\ref{tab:setup}. All three modify existing nodes; none adds nodes to a
graph. None of the attacks is adapted to the detectors.

\subsection{Simulation Results}
% Source: results/analysis/paper_graph_20260930 (+ _logreg); see the
% provenance note at the top of this file. One table answers all three
% questions; each subsubsection below reads one column group of it.

\begin{table*}[t]
\caption{Detection and defense results for a GCN target model, averaged over
five folds. CA and ASR: clean accuracy and attack success rate of the poisoned
model without defense. FPR and TPR (\%): clean and triggered test graphs
flagged. AUC: ROC-AUC; the statistical detector ranks graphs by its smallest
Holm-adjusted $p$-value. ASR after: attack success rate when flagged graphs are
rejected. The MLP threshold is calibrated to the clean-rejection rate of the
statistical detector; logistic regression (LR) is the linear reference.}
\label{tab:results}
\centering
\tabstyle
\begin{tabular}{llcc cccc c cccc}
\toprule
         &        &       &       &
\multicolumn{4}{c}{Statistical detector} &
LR &
\multicolumn{4}{c}{MLP detector} \\
\cmidrule(lr){5-8}\cmidrule(lr){9-9}\cmidrule(lr){10-13}
Dataset  & Attack & CA & ASR &
FPR & TPR & AUC & ASR after &
AUC &
AUC & FPR & TPR & ASR after \\
\midrule
AIDS     & SBA   & 0.890 & 0.086 &   5.8 &  15.6 & 0.615 & 0.073 & 0.869 & 0.910 &   5.1 &  67.6 & 0.021 \\
         & Motif & 0.896 & 0.090 &   5.8 &  32.8 & 0.733 & 0.065 & 0.977 & 0.993 &   4.9 &  99.4 & 0.000 \\
         & GTA   & 0.903 & 0.999 &   5.8 &  76.0 & 0.939 & 0.240 & 0.995 & 0.999 &   6.3 &  99.9 & 0.001 \\
\midrule
MUTAG    & SBA   & 0.745 & 0.152 &   7.5 &  68.8 & 0.886 & 0.064 & 0.985 & 0.983 &   8.0 &  96.0 & 0.000 \\
         & Motif & 0.750 & 0.152 &   7.5 &  95.2 & 0.985 & 0.000 & 1.000 & 1.000 &   8.0 & 100.0 & 0.000 \\
         & GTA   & 0.756 & 1.000 &   7.5 &  96.8 & 0.986 & 0.032 & 1.000 & 1.000 &   6.9 & 100.0 & 0.000 \\
\midrule
NCI1     & SBA   & 0.680 & 0.330 &   9.2 &  21.0 & 0.621 & 0.254 & 0.868 & 0.913 &   8.7 &  76.3 & 0.083 \\
         & Motif & 0.681 & 0.518 &   9.2 &  30.8 & 0.681 & 0.348 & 0.935 & 0.990 &   9.0 &  99.2 & 0.004 \\
         & GTA   & 0.672 & 1.000 &   9.2 &  42.5 & 0.747 & 0.575 & 0.977 & 0.996 &   9.7 &  99.7 & 0.003 \\
\midrule
PROTEINS & SBA   & 0.697 & 0.518 &   3.1 &   7.6 & 0.553 & 0.489 & 0.830 & 0.883 &   3.2 &  40.9 & 0.300 \\
         & Motif & 0.704 & 0.484 &   3.1 &  11.8 & 0.553 & 0.449 & 0.761 & 0.844 &   3.1 &  34.2 & 0.307 \\
         & GTA   & 0.724 & 0.951 &   3.1 &  11.3 & 0.603 & 0.844 & 0.842 & 0.881 &   3.2 &  49.8 & 0.462 \\
\bottomrule
\end{tabular}
\end{table*}

\subsubsection{Statistical Detection}
The statistical detector asks whether a triggered graph departs from the
clean reference at some coordinate by more than clean graphs do. Every eligible
triggered test graph is a clean test graph edited only by the trigger, so a
triggered rejection rate above the clean rate shows that the trigger's spectral
fingerprint lies outside the clean variation. Without multiple-testing
correction, the nominal coordinatewise level flags $30.8\%$--$64.3\%$ of clean
test graphs, because every graph is tested at many coordinates. With Holm
correction, the clean rejection falls to $3.1\%$--$9.2\%$
(Table~\ref{tab:results}). At this cost, the detector flags
$68.8\%$--$96.8\%$ of triggered MUTAG graphs, and GTA yields the highest
triggered rejection on AIDS and NCI1, at $76.0\%$ and $42.5\%$, respectively.
Under SBA and Motif-Backdoor on AIDS and NCI1, it flags only
$15.6\%$--$32.8\%$ of triggered graphs, which is $2.3$ to $5.7$ times its clean
rejection. On PROTEINS, the triggered rejection of $7.6\%$--$11.8\%$ is $2.5$ to
$3.9$ times the clean rejection of $3.1\%$, and the ROC-AUC of its ranking score
is only $0.553$--$0.603$. Trigger insertion thus moves the descriptor beyond
the natural clean variation for most triggered graphs only on MUTAG and under
GTA on AIDS; in the remaining settings, most triggered graphs have no single
coordinate that deviates enough to survive the correction.

\subsubsection{MLP Detection Performance}
The MLP reaches ROC-AUC values of $0.844$--$1.000$ and, at its calibrated
threshold, detects $34.2\%$--$100.0\%$ of triggered graphs at an FPR of
$3.1\%$--$9.7\%$. Its ROC-AUC exceeds that of logistic regression in $9$ of
the $12$ settings, with the largest gains under Motif-Backdoor on PROTEINS
($0.761$ to $0.844$) and NCI1 ($0.935$ to $0.990$). This gain
indicates that part of the discriminative information lies in nonlinear
combinations of coordinates. On MUTAG, both models separate the graphs almost
perfectly, and logistic regression is slightly higher under SBA; the MLP thus
does not uniformly dominate the linear reference.

PROTEINS remains the most difficult dataset. At an FPR of $3.1\%$--$3.2\%$, the
MLP detects $34.2\%$--$49.8\%$ of triggered graphs, whereas the statistical
detector, at the same clean cost, detects $7.6\%$--$11.8\%$. Part of this gap
comes from supervision, since the statistical detector uses no triggered
examples: logistic regression already reaches ROC-AUC $0.761$--$0.842$, against
$0.553$--$0.603$ for the statistical detector, and the MLP adds $0.04$--$0.08$.
Because no single coordinate separates most triggered PROTEINS graphs from clean
ones, this indicates that weak deviations spread across the descriptor remain
jointly discriminative once labelled examples are available.

\subsubsection{Analyzing Overall Defense Capability}
The ASR columns of Table~\ref{tab:results} show the effect of rejecting
flagged graphs, and the reduction follows the separation reported above.
Under GTA, which succeeds on almost every eligible graph, the MLP lowers ASR
from $0.999$ to $0.001$ on AIDS, from $1.000$ to $0.000$ on MUTAG, and from
$1.000$ to $0.003$ on NCI1, whereas the statistical detector leaves $0.240$,
$0.032$, and $0.575$, respectively. On PROTEINS under GTA, rejection lowers ASR
from $0.951$ to $0.844$ with the statistical detector and to $0.462$ with the
MLP, so substantial backdoor behavior remains. The same holds for SBA and
Motif-Backdoor on PROTEINS, with post-detection ASR between $0.300$ and
$0.489$. SBA and Motif-Backdoor succeed on few AIDS and MUTAG graphs even
without a defense (ASR $0.086$--$0.152$), so the absolute reductions there are
small.

The MLP yields a post-detection ASR no higher than that of the statistical
detector in every setting; the two are equal only on MUTAG under
Motif-Backdoor, where both reach $0.000$. This comparison is descriptive: the
five folds are not independent experiments, and the FPRs of the two detectors
differ by less than one percentage point. Spectral rejection therefore suppresses
the backdoor when the triggered descriptor leaves the range of clean
variation, and its benefit shrinks as the clean and triggered distributions
overlap.

\section{Conclusion}
This letter examined whether structural backdoor triggers in graph
classification remain imperceptible in the spectral domain. On the molecular
benchmarks, a spectrum-only MLP separates triggered from clean graphs with
ROC-AUC between $0.91$ and $1.00$ and, when used to reject inputs, removes
nearly all success of the learned GTA trigger on AIDS, MUTAG, and NCI1 while
rejecting $6.3\%$--$9.7\%$ of clean graphs. The clean-reference test is
weaker: it flags most triggered graphs only on MUTAG and under GTA on AIDS. On
PROTEINS, both detectors are less effective and substantial attack success
remains after rejection. Three limitations bound these findings. The MLP is
trained on triggered graphs of the attack it detects, so its results bound what
the spectrum reveals rather than describe a defense against unseen attacks, and
transfer across attacks is not addressed. The attacks are not adapted to the
detectors, and only a GCN target model and structural triggers are considered.
No non-spectral baseline is compared, so the detectors may respond to any
structural change that the spectrum reflects, such as a changed edge count.
Detecting triggers without labelled examples, and against adaptive attackers,
remains open.

% Optional -- IEEE prefers funding acknowledgments in the first-page
% footnote (\thanks) rather than here.
% \section*{Acknowledgment}

% =====================================================================
% References. IEEEtran.bst is fetched by tectonic and ships with both
% TeX Live and Overleaf, so it is not vendored in this repo.
%
% main.bib holds the curated source records for this letter. Verify each
% record against the publisher's own export before submission.
% =====================================================================
\newpage
\bibliographystyle{IEEEtran}
\bibliography{main}
\newpage
\appendices

\section{Training Settings}
\label{app:training}
The five folds are stratified, and within each fold a nested stratified
split reserves $10\%$ of the full dataset for validation. The split seed is
$12345$ and the model seed is $0$. The reported results use a GCN target
model at trigger size $k=4$; runs with GraphSAGE~\cite{graphsage} and
GAT~\cite{gat} target models and at other trigger sizes are not included in
this letter.

The graph-attack task file configures SBA with an Erd\H{o}s--R\'enyi pattern
and Motif-Backdoor with the M44 motif. The run records store the shared budget
and target settings but not the remaining attack parameters, which are
implementation defaults: SBA edge probability $0.5$, Motif-Backdoor candidate
pool of $10$ nodes, and, for GTA, two-hop adjacency and aggregated node
features as generator input, topology threshold $0.5$, feature threshold $0$,
and generator learning rate $0.01$. For a graph with fewer than $k$ nodes, SBA
inserts the pattern induced on the available nodes if there are at least two;
otherwise, and for Motif and GTA, the graph is left unchanged.

\section{Benchmark Structure}
\label{app:structure}
Table~\ref{tab:dataset_stats_graph} summarizes the benchmarks. PROTEINS has
larger graphs and a higher mean degree than the molecular benchmarks. A fixed
trigger budget can therefore represent a smaller relative structural change;
this is a possible contributor to its weaker detection, not a causal
conclusion established by these comparisons.

\begin{table}[t]
\caption{Graph-classification benchmarks. Averages are over all graphs in the
dataset; $|\calE|$ counts undirected edges, and average degree is the mean over
graphs of each graph's $2|\calE|/n$.}
\label{tab:dataset_stats_graph}
\centering
\tabstyle
\begin{tabular}{@{}lrrrrrr@{}}
\toprule
Dataset & Graphs & Avg.\ $n$ & Avg.\ $|\calE|$ & Avg.\ deg. & Feat. & Classes \\
\midrule
AIDS     & 2000 & 15.69 & 16.20 & 2.01 & 38 & 2 \\
MUTAG    &  188 & 17.93 & 19.79 & 2.19 &  7 & 2 \\
NCI1     & 4110 & 29.87 & 32.30 & 2.16 & 37 & 2 \\
PROTEINS & 1113 & 39.06 & 72.82 & 3.73 &  3 & 2 \\
\bottomrule
\end{tabular}
\end{table}

\end{document}
